% !TEX root = ../main.tex
\section{Results}

We report the reproduced findings of the full experiment grid (18 landbases $\times$ 4 discount rates $\times$ 6 harvest-flow policies; 432 sequential-replanning simulations). Per the framing of this paper, these reproduce and refine \citet{daugherty1991credibility}'s findings on the real Umpqua FORPLAN case study; they are presented with their fidelity caveats in Section 5. Every number below is regenerated by the repository's tracked experiment entry point; per-cell simulation records are in the repository.

\textbf{The harvest-flow constraint is the operative link.} Dynamic inconsistency is pervasive under the harvest-flow policies but largely absent without one. Across the grid, 73\% of the flow-constrained cells (NDY, bounded decline, and bounded deviation) exhibit inconsistency by the divergence metric (equation~\eqref{eq:metric}), against 38\% of the no-harvest-flow (NHF) control cells---and the NHF divergence is concentrated at the lowest discount rates (0--2\%), where the objective is nearly indifferent to timing; at 4--6\% the no-flow control is consistent (Section 5). In all flow-constrained cells the first-period decision is consistent---the divergence is in the plan's tail, exactly as the theory predicts. Where a model links periods only through a harvest-flow constraint, its projected harvest trajectory is a property of a plan that sequential re-optimization will not follow.

\textbf{The declining non-declining yield.} On the all-mature landbase 1, the divergence between announcement and delivery is stark and monotone (Figure~\ref{fig:ndy}). The open-loop non-declining-yield plan projects a smooth, credible-looking level harvest flow; the realized replanned trajectory is revised steadily downward as each successive planner inherits the depleted forest and re-solves. This is the ``declining non-declining yield'' of \citet[pp.~331--332]{gunn2007models}: the promised non-declining flow declines precisely because each successive planner declines to honour it.

\begin{figure}[t]
\centering
\includegraphics[width=\textwidth]{figures/fig_declining_ndy.pdf}
\caption{The declining non-declining yield on the all-mature landbase 1 under the non-declining-yield (NDY) harvest-flow policy at a 4\% discount rate. The open-loop plan projects a level, non-declining harvest flow; the realized sequential-replanning trajectory is revised steadily downward. Full provenance and the per-cell simulation records are in the project repository.}
\label{fig:ndy}
\end{figure}

\textbf{Occurrence at every discount rate; magnitude is rate-dependent.} Inconsistency occurs at every discount rate in the grid (the flow-constrained occurrence rate is 100\% at 0\% and 47--86\% at 2--6\%), consistent with the thesis's analytical result that, because the discount factor is uniform across periods, it drops out of the consistency requirements \citep[ch.~3, p.~58--59]{daugherty1991credibility}---the phenomenon does not require a particular rate. The \emph{magnitude} is rate-dependent: on landbase 1 under NDY, the mean relative divergence and the total-volume shortfall are markedly larger at a 0\% rate (where the objective is indifferent to timing) than at 2--6\%. We return to this nuance in the Discussion, because it refines---rather than contradicts---the naive account of the phenomenon.

\textbf{Effect of the harvest-flow policy.} The strictness and form of the flow constraint shape the inconsistency. The non-declining-yield policy---permitting no decrease---produces the most frequent divergence (occurrence 86\%, the highest of the flow policies), and the magnitude of the divergence varies by policy and landbase (Table 1). The no-flow control is the least inconsistent throughout. This confirms the flow constraint as the operative between-period link through which current decisions distort future subproblems.

\textbf{Role of disequilibrium structure.} The inconsistency occurs across all eighteen initial forest conditions (flow-constrained occurrence 50--100\% by landbase), and is most pronounced in magnitude on the disequilibrium structures: the all-mature and area-control-derived landbases, which place a pulse of immediately harvestable, financially over-mature volume in front of a regulated future that cannot sustain it. The CM-CE ecoclass is negatively valued---its harvest cost exceeds its timber value---and such strata enter the open-loop solutions as \emph{inconsistent variables}: area allocations whose only function is to relax the binding flow constraint in the plan and that the re-solving planner subsequently abandons.

\textbf{Table 1.} Experiment-grid summary: factors, levels, occurrence of dynamic inconsistency (share of cells exceeding the 5\% divergence tolerance), and mean magnitude (symmetric relative deviation, equation~\eqref{eq:metric}), over the full grid (432 cells). Full per-cell simulation records are in the project repository.

\begin{table}[h]
\centering
\small
\begin{tabular}{@{}llcc@{}}
\toprule
Factor & Levels & Occurrence & Mean magnitude \\
\midrule
Harvest-flow policy & NDY, $-10\%$, $-20\%$, $\pm10\%$, $\pm20\%$ & 58--86\% & 0.10--0.13 \\
\quad (control) & NHF (no harvest flow) & 38\% & 0.13 \\
Discount rate & 0\%, 2\%, 4\%, 6\% & 100\% / 47--86\% & rate-dependent \\
Initial forest condition & 18 landbases (Table 5.5) & 50--100\% & 0.08--0.23 \\
\bottomrule
\end{tabular}
\end{table}
